\input{cyclic/continuation_20261026/injectivity45.tex}
\begin{proposition}\label{prop:nonic-no45}
An elliptic curve over a number field of degree $9$ has no point of order $45$.
\end{proposition}
\begin{proof}
Put $X=X_1(45)$. It has genus $41$, rank-zero rational Jacobian
by \cite[Theorem~3.1(2)]{DEHMZ}, and rational gonality at least $10$
by \cite[Table~1]{NV}. At $2$ its noncuspidal closed-place counts
through degree $9$ are
\[
 (0,0,0,0,0,0,0,60,60).
\]
Independent Frobenius similarity-class counts agree.
Each of degrees $8,9$ has five diamond orbits.
For all five degree-$8$ representatives the immutable complete
finite-field computation finds a nonprincipal permitted $A_q$ image.
The only new complement has degree $1$ and is a rational cusp, killed
by every tested $A_q$; thus the same tests apply in degree $9$.
For all five degree-$9$ representatives, fresh full-source
Riemann--Roch tests give a nonprincipal $A_7$ image.
An untouched native $7$-isogeny calculation independently reproduces
the degree-$72$ Hecke divisor of one representative.
The complete fresh run finishes with exit code $0$.
Thus no noncuspidal reduction can occur.

The finite cusp signature is $1^{12},2^4,4^5,6^2,12$.
Put $U=(\mathcal C-1)H^0(X_{\mathbb F_2},\Omega^1)$, of dimension $39$.
The degree-independent local restrictions of weighted multiplicity
at most $8$ are retained in their single verified global basis.
The independent binary calculation enumerates $377348$ degree-$9$
cuspidal divisors covered by these restrictions. Exactly eight have
rank $8$; every other restriction has rank $9$.
A fresh native differential-space kernel at one rational cusp gives
rank $9$ at multiplicity $9$. Diamond transitivity gives the same
conclusion for the twelve pure divisors $9c$.
Lemma~\ref{lem:nonic51-quotient-Cartier} therefore excludes every
entirely cuspidal reduction except the eight exported configurations.
No bases from different geometric exports are combined.

Each of these eight configurations is supported only on rational
and quadratic residue cusps. The global cusp degrees and their
Frobenius-$2$ degrees, by \cite[Lemma~2.7]{DEHMZ}, are
\[
 (12,12),\ (8,4),\ (6,6)^2,\ (4,4)^3,
 \ (2,2)^4,\ (1,1)^{12}.
\]
Consequently every exceptional configuration lifts to a rational
effective cuspidal divisor $C$ of degree $9$.
If a rational noncuspidal divisor $D$ had this reduction,
$[D-C]\in J_1(45)(\mathbb Q)$ would reduce to zero.
Lemma~\ref{lem:nonic45-injective-two} gives $D\sim C$.
Since the gonality is greater than $9$, the two effective divisors
are equal, a contradiction.
This also excludes the eight remaining configurations and proves
the proposition.
\end{proof}
